\chapter{Charakterisierungen und Grundregeln}
\label{b14guide:ch:characterizations}

\section{Universelle Charakterisierungen}

\begin{theorem}[Charakterisierung eines Kandidaten]
\label{b14guide:thm:candidate}
Für $m,s\in P$ gilt, ohne vorherige Existenzannahme,
\[
\begin{aligned}
 m=x\wedge y
 &\quad\Longleftrightarrow\quad
 \forall z\in P\;\bigl(z\leq m\Longleftrightarrow
                       (z\leq x\land z\leq y)\bigr),\\
 s=x\vee y
 &\quad\Longleftrightarrow\quad
 \forall z\in P\;\bigl(s\leq z\Longleftrightarrow
                       (x\leq z\land y\leq z)\bigr).
\end{aligned}
\]
Die rechte Seite sichert jeweils sowohl die Existenz als auch
den angegebenen Wert.
\end{theorem}

\begin{theorem}[Auswertungs- und Vergleichsregeln]
\label{b14guide:thm:universal}
Ist $x\wedge y$ vorhanden, so gelten
\[
 x\wedge y\leq x,\qquad x\wedge y\leq y,\qquad
 z\leq x\wedge y\Longleftrightarrow(z\leq x\land z\leq y).
\]
Ist $x\vee y$ vorhanden, so gelten
\[
 x\leq x\vee y,\qquad y\leq x\vee y,\qquad
 x\vee y\leq z\Longleftrightarrow(x\leq z\land y\leq z).
\]
\end{theorem}

\begin{theorem}[Paarwerte und Hauptabschnitte]
\label{b14guide:thm:sections}
Schreiben wir $\mathord\downarrow a=\{z\in P\mid z\leq a\}$
und $\mathord\uparrow a=\{z\in P\mid a\leq z\}$, so gilt
\[
\begin{aligned}
 m=x\wedge y
 &\Longleftrightarrow
 \mathord\downarrow m
   =\mathord\downarrow x\cap\mathord\downarrow y,\\
 s=x\vee y
 &\Longleftrightarrow
 \mathord\uparrow s
   =\mathord\uparrow x\cap\mathord\uparrow y.
\end{aligned}
\]
Ein Paarinfimum existiert also genau dann, wenn der erste
Durchschnitt ein Hauptabschnitt dieser Form ist; für das
Paarsupremum gilt die zweite Aussage entsprechend.
\end{theorem}

\section{Vergleichbare und gleiche Operanden}

\begin{theorem}[Vergleichbare Operanden]
\label{b14guide:thm:comparable}
Aus $x\leq y$ folgt die Existenz beider Paarwerte, und es gilt
\[
 x\wedge y=x,\qquad x\vee y=y.
\]
Insbesondere existieren für jedes $x\in P$ die Werte
\[
 x\wedge x=x,\qquad x\vee x=x
 \qquad\text{(Idempotenz).}
\]
\end{theorem}

\begin{theorem}[Erkennen der Ordnung an einem Paarwert]
\label{b14guide:thm:order}
Für jedes $x,y\in P$ sind äquivalent:
\[
 x\leq y,\qquad
 x\wedge y\text{ existiert und ist gleich }x,\qquad
 x\vee y\text{ existiert und ist gleich }y.
\]
Jede der letzten beiden Aussagen erzwingt somit auch die
Existenz des jeweils anderen Paarwertes.
\end{theorem}

\begin{theorem}[Vertauschen der Operanden]
\label{b14guide:thm:commutative}
Es gilt
\[
 \InfEx(x,y)\Longleftrightarrow\InfEx(y,x),\qquad
 \SupEx(x,y)\Longleftrightarrow\SupEx(y,x).
\]
Im jeweiligen Existenzfall gilt die Kommutativität
\[
 x\wedge y=y\wedge x,\qquad x\vee y=y\vee x.
\]
\end{theorem}

\begin{theorem}[Strikte Vergleiche und Unvergleichbarkeit]
\label{b14guide:thm:strict}
Es sei $a<b$ die Abkürzung für $a\leq b$ und $a\ne b$.
Ist $x\wedge y$ vorhanden, so gilt
\[
 x\wedge y<x\Longleftrightarrow x\not\leq y,\qquad
 x\wedge y<y\Longleftrightarrow y\not\leq x.
\]
Ist $x\vee y$ vorhanden, so gilt
\[
 x<x\vee y\Longleftrightarrow y\not\leq x,\qquad
 y<x\vee y\Longleftrightarrow x\not\leq y.
\]
Bei unvergleichbaren Operanden liegt jeder vorhandene
Paarinfimumswert strikt unter beiden, jeder vorhandene
Paarsupremumswert strikt über beiden Operanden.
\end{theorem}

\begin{theorem}[Zusammenfallen von Paarinfimum und Paarsupremum]
\label{b14guide:thm:coincidence}
Sind beide Paarwerte vorhanden, so gilt
\[
 x\wedge y\leq x\leq x\vee y,\qquad
 x\wedge y\leq y\leq x\vee y,
\]
und
\[
 x\wedge y=x\vee y\Longleftrightarrow x=y.
\]
Für $x\ne y$ folgt daher $x\wedge y<x\vee y$.
\end{theorem}